% ===================================================================
% 线性代数（数学三）—— 插入到主文档「十一、线性代数」处
% 依赖主文档已有宏：\secbox  \formbox  \bx  \cond  \hd
% ===================================================================

% -------------------------------------------------------------------
\begin{secbox}[十一、线性代数 · \textcircled{1} 行列式]

\begin{formbox}
\bx{
|A|=\sum_{j=1}^{n}a_{1j}A_{1j}
=\sum_{i=1}^{n}a_{i1}A_{i1}
}
\hfill（按任一行或任一列展开）

\bx{
\sum_{j=1}^{n}a_{ij}A_{kj}=\delta_{ik}\,|A|
}
\hfill\note{（$i,k\in\{1,\dots,n\}$，即 $i=k$ 时等于 $|A|$，$i\ne k$ 时等于 $0$）}
\end{formbox}

\begin{formbox}
\renewcommand{\arraystretch}{1.5}
\begin{tabular}{@{}p{7.4cm}p{8.6cm}@{}}
$|A|=|A^{T}|$
  & 转置行列式不变 \\[4pt]
$|kA|=k^{n}|A|$
  & 数乘：$n$ 阶乘 $k^{n}$ \\[4pt]
$|AB|=|A||B|$
  & 乘积行列式 \\[4pt]
$|A+B|\ne|A|+|B|$
  & \textbf{不可}像行列式那样拆开相加 \\[4pt]
交换任意两行（列）$\Rightarrow|A|$ 变号
  & $k$ 次交换则乘 $(-1)^{k}$ \\[4pt]
某行（列）$\times k\Rightarrow|A|\times k$
  & 某行（列）$\div k\Rightarrow|A|\div k$ \\[4pt]
某行（列）$+k\times$另一行（列）$\Rightarrow|A|$ 不变
  & 连续作两次 $\Rightarrow$ 不变
\end{tabular}
\end{formbox}

\begin{formbox}
\renewcommand{\arraystretch}{1.6}
\[
\text{三角行列式：}\quad
\left|\begin{matrix}a_{11}& & \\ &\ddots& \\ &&a_{nn}\end{matrix}\right|
=\prod_{i=1}^{n}a_{ii}
\qquad\text{（上、下三角同）}
\]
\[
\text{范德蒙德行列式：}\quad
\left|\begin{matrix}
1&1&\cdots&1\\
x_{1}&x_{2}&\cdots&x_{n}\\
x_{1}^{2}&x_{2}^{2}&\cdots&x_{n}^{2}\\
\vdots&\vdots&&\vdots\\
x_{1}^{n-1}&x_{2}^{n-1}&\cdots&x_{n}^{n-1}
\end{matrix}\right|
=\prod_{1\le i<j\le n}\left(x_{j}-x_{i}\right)
\]
\[
\text{分块三角行列式：}\quad
\begin{vmatrix}A&B\\ O&C\end{vmatrix}=|A||C|,
\qquad
\begin{vmatrix}A&O\\ C&D\end{vmatrix}=|A||D|
\]
\[
\text{反对称矩阵（}\ n\ \text{为奇数）：}\quad
A^{T}=-A\ \Rightarrow\ |A|=0
\qquad
\text{（偶数阶一般不为零）}
\]
\[
\text{含两行（列）相同或成比例：}\quad |A|=0
\qquad
\text{含零行（列）：}\quad |A|=0
\]
\end{formbox}
\end{secbox}

% -------------------------------------------------------------------
\begin{secbox}[十二、线性代数 · \textcircled{2} 矩阵]

\begin{formbox}
\renewcommand{\arraystretch}{1.5}
\begin{tabular}{@{}p{7.4cm}p{8.6cm}@{}}
$(AB)^{T}=B^{T}A^{T}$
  & 转置反序 \\[3pt]
$(kA)\pm lB=k(A\pm lB)$，$AB+AC=A(B+C)$
  & 线性运算 \\[3pt]
$(AB)C=A(BC)$
  & 结合律成立、乘法不满足交换律 \\[3pt]
$kA\cdot lB=kl\,(AB)$
  & 数乘可提
\end{tabular}

\medskip
\hd{逆矩阵（$A$ 为 $n$ 阶方阵，$|A|\ne0$）}
\[
(A^{-1})^{-1}=A,\qquad (A^{T})^{-1}=(A^{-1})^{T},
\qquad (kA)^{-1}=\frac{1}{k}A^{-1}
\]
\[
\bx{(AB)^{-1}=B^{-1}A^{-1}}\quad\text{（逆序！高频错点）}
\]
\end{formbox}

\begin{formbox}
$A$ 可逆的等价结论（任一条成立即可断定可逆）：
\[
\bx{
|A|\ne0\iff r(A)=n\iff Ax=\boldsymbol0\ \text{只有零解}
}
\]
\[
\bx{
A\ \text{的列向量组线性无关}\iff A\ \text{的行向量组线性无关}
}
\]
\[
\bx{
\lambda_{1},\dots,\lambda_{n}\ \text{均非零}
\iff r(\lambda E-A)=n,\ \forall\lambda
}
\]
\end{formbox}

\begin{formbox}
\hd{伴随矩阵}
\[
A^{*}=\left(a_{ij}^{*}\right),\qquad a_{ij}^{*}=A_{ji}
\qquad\text{（余子式转置；代数余子式 $A_{ij}$ 不转置）}
\]
\[
\bx{AA^{*}=A^{*}A=|A|E},
\qquad
A^{-1}=\frac{1}{|A|}A^{*}\ \ (|A|\ne0),
\qquad
|A^{*}|=|A|^{\,n-1}
\]
\[
(A^{*})^{*}=|A|^{\,n-2}A,
\qquad
(kA)^{*}=k^{\,n-1}A^{*},
\qquad
\bx{(AB)^{*}=B^{*}A^{*}}
\]
\end{formbox}

\begin{formbox}
\hd{矩阵的秩}
\[
\begin{aligned}
r(A)&=r(A^{T})\\
r(A+B)&\le r(A)+r(B)\\
\bx{r(AB)}&\le\min\{r(A),\,r(B)\}\\
r(AB)&\le r(A),\qquad r(AB)\le r(B)\\
A\ \text{可逆}&\ \Rightarrow\ r(AB)=r(BA)=r(B)\\
AB=\boldsymbol 0&\ \Rightarrow\ r(A)+r(B)\le n
\qquad\text{（$A$ 为 $m\times n$ 矩阵）}\\
r(A+E)&\le r(A)+1,\qquad r(A\cdot B)\le r(A)\cdot r(B)
\end{aligned}
\]
\[
\bx{r(A)=n\iff |A|\ne0\iff A\ \text{可逆}\iff Ax=\boldsymbol 0\ \text{只有零解}}
\]
\end{formbox}

\begin{formbox}
\hd{初等变换、初等矩阵与等价}
\[
\begin{gathered}
\text{初等变换：}\quad
\text{交换两行（列）}\ \big|\ \text{某行（列）}\times k\ (k\ne0)\ \big|\ \text{某行（列）}+k\times\text{另一行（列）}\\[4pt]
E\ \text{可逆},\qquad r(EA)=r(A),\qquad E_{i}\cdots E_{k}A=E_{i}\cdots E_{k}B\iff A\ \text{与}\ B\ \text{等价}\\
A\sim B\iff \exists\,P,Q\ \text{可逆},\ PAQ=B
\iff \text{同型且}\ r(A)=r(B)
\end{gathered}
\]
\[
\bx{\text{求秩（通用法）：}\quad A\xrightarrow{\text{初等行变换}}\text{行最简形}\ \Rightarrow\ r(A)=\text{非零行数}=\text{非零列数}}
\]
\[
\text{求逆（初等行变换法）：}\quad
\left(\begin{array}{c|c}A&E\end{array}\right)
\xrightarrow{\text{仅行变换}}
\left(\begin{array}{c|c}E&A^{-1}\end{array}\right)
\]
\end{formbox}

\begin{formbox}
\hd{分块矩阵}
\[
\begin{vmatrix}A&B\\ O&C\end{vmatrix}
=\begin{vmatrix}A&O\\ C&D\end{vmatrix}=|A||C|,
\qquad
\begin{pmatrix}A&B\\ O&C\end{pmatrix}^{-1}
=\begin{pmatrix}A^{-1}&-A^{-1}BC^{-1}\\ O&C^{-1}\end{pmatrix}
\]
\[
\begin{pmatrix}A&O\\ O&D\end{pmatrix}^{-1}
=\begin{pmatrix}A^{-1}&O\\ O&D^{-1}\end{pmatrix},
\qquad
\begin{pmatrix}A&B\\ C&D\end{pmatrix}^{T}
=\begin{pmatrix}A^{T}&C^{T}\\ B^{T}&D^{T}\end{pmatrix}
\]
\[
A\ \text{对称}\iff
\begin{pmatrix}A&B\\ C&A\end{pmatrix}\ \text{对称}
\iff A^{T}=A,\ B^{T}=C
\]
\end{formbox}
\end{secbox}

% -------------------------------------------------------------------
\begin{secbox}[十三、线性代数 · \textcircled{3} 向量]

\begin{formbox}
$\boldsymbol\beta=k_{1}\boldsymbol\alpha_{1}+\cdots+k_{s}\boldsymbol\alpha_{s}$
称 $\boldsymbol\beta$ 可由 $\boldsymbol\alpha_{1},\dots,\boldsymbol\alpha_{s}$ \textbf{线性表示}。
\begin{align*}
\bx{\boldsymbol\alpha_{1},\dots,\boldsymbol\alpha_{s}\ \text{线性相关}}
&\iff \exists\ k_{1},\dots,k_{s}\ \text{不全为}\ 0,\ \sum_{i=1}^{s}k_{i}\boldsymbol\alpha_{i}=\boldsymbol 0\\
&\iff r(\boldsymbol\alpha_{1},\dots,\boldsymbol\alpha_{s})<s\\
\bx{\boldsymbol\alpha_{1},\dots,\boldsymbol\alpha_{s}\ \text{线性无关}}
&\iff \sum_{i=1}^{s}k_{i}\boldsymbol\alpha_{i}=\boldsymbol 0\ \Rightarrow\ k_{1}=\cdots=k_{s}=0\\
&\iff r(\boldsymbol\alpha_{1},\dots,\boldsymbol\alpha_{s})=s
\end{align*}
\[
s>n\ \Rightarrow\ \boldsymbol\alpha_{1},\dots,\boldsymbol\alpha_{s}\ \text{必线性相关}
\qquad
\text{（}\ n\ \text{维空间中最多}\ n\ \text{个无关向量）}
\]
\[
\text{整体无关}\ \Rightarrow\ \text{部分必无关};
\qquad \text{部分相关}\ \Rightarrow\ \text{整体必相关};
\qquad \text{反之均不成立}
\]
\[
n\ \text{个}\ n\ \text{维向量无关}\iff |\boldsymbol\alpha_{1},\dots,\boldsymbol\alpha_{n}|\ne0
\]
\[
\boldsymbol\alpha_{1},\dots,\boldsymbol\alpha_{s}\ \text{无关},\quad
\boldsymbol\beta_{1},\dots,\boldsymbol\beta_{t}\ \text{无关},\quad
s+t>n\ \Rightarrow\ \text{两组必线性相关}
\]
\end{formbox}

\begin{formbox}
\hd{极大线性无关组、秩与等价}
\[
r(\boldsymbol\alpha_{1},\dots,\boldsymbol\alpha_{s})
=r(\text{以它们为列的矩阵})
=r(\text{行向量组})=r(\text{列向量组})
\]
\[
r(A)=\text{极大无关组的个数}=\text{非零子式的最高阶数}
\]
\[
\begin{aligned}
&\text{无关组}\ (\mathrm I)\ \text{可由相关组}\ (\mathrm{II})\ \text{线性表出}
\ \Rightarrow\ r(\mathrm I)\le r(\mathrm{II})\\[3pt]
&\text{若反之亦可表出}\ \Rightarrow\ r(\mathrm I)=r(\mathrm{II})\ \text{（此时两组等价）}\\[3pt]
&\text{等价}\ \Rightarrow\ \text{秩相等；}\qquad
\text{秩相等}\ \not\Rightarrow\ \text{等价}\ \text{（反向不成立）}
\end{aligned}
\]
\end{formbox}

\begin{formbox}
\hd{内积与正交规范化（施密特 Schmidt 法）}
\[
(\boldsymbol\alpha,\boldsymbol\beta)
=\boldsymbol\alpha^{T}\boldsymbol\beta=\sum_{i=1}^{n}\alpha_{i}\beta_{i}
\qquad
\boldsymbol\alpha\perp\boldsymbol\beta\iff(\boldsymbol\alpha,\boldsymbol\beta)=0
\]
\[
\boldsymbol\xi_{1}=\boldsymbol\alpha_{1},\qquad
\boldsymbol\xi_{k}=\boldsymbol\alpha_{k}
-\sum_{i=1}^{k-1}\frac{(\boldsymbol\alpha_{k},\boldsymbol\xi_{i})}{(\boldsymbol\xi_{i},\boldsymbol\xi_{i})}\boldsymbol\xi_{i},
\quad
\bx{\boldsymbol\eta_{k}=\frac{\boldsymbol\xi_{k}}{\lVert\boldsymbol\xi_{k}\rVert}}
\qquad(k=2,3,\dots,s)
\]
\[
\lVert\boldsymbol\alpha\rVert=\sqrt{(\boldsymbol\alpha,\boldsymbol\alpha)},
\qquad
(\boldsymbol\alpha,\boldsymbol\beta)=\lVert\boldsymbol\alpha\rVert\lVert\boldsymbol\beta\rVert\cos\theta
\]
\end{formbox}
\end{secbox}

% -------------------------------------------------------------------
\begin{secbox}[十四、线性代数 · \textcircled{4} 线性方程组]

\begin{formbox}
设 $m$ 个方程、$n$ 个未知数：
\[
\begin{array}{l}
a_{11}x_{1}+\cdots+a_{1n}x_{n}=b_{1}\\
\ \vdots\qquad\vdots\qquad\ \vdots\\
a_{m1}x_{1}+\cdots+a_{mn}x_{n}=b_{m}
\end{array}
\quad\Longleftrightarrow\quad
Ax=\boldsymbol b,\qquad A\ \text{为}\ m\times n\ \text{矩阵}
\]
\[
\text{克莱姆（Cramer）法则：}\quad |A|\ne0\ \Rightarrow\ x_{i}=\frac{|A_{i}|}{|A|}
\qquad\text{（$A_{i}$ 为 $A$ 中第 $i$ 列换成 $\boldsymbol b$ 所得）}
\]
\end{formbox}

\begin{formbox}
\hd{齐次方程组 $Ax=\boldsymbol 0$ 的结论}
\[
\bx{
|A|=0\iff r(A)<n\iff Ax=\boldsymbol 0\ \text{有非零解}
\iff A\ \text{的列向量组线性相关}
}
\]
\[
n-r(A)\ \text{个线性无关的解向量构成\textbf{基础解系}}
\qquad
\bx{\text{通解}\ x=k_{1}\boldsymbol\eta_{1}+\cdots+k_{n-r(A)}\boldsymbol\eta_{n-r(A)}}
\]
\end{formbox}

\begin{formbox}
\hd{非齐次方程组 $Ax=\boldsymbol b$ 的解的判定}
\[
\bxe{
\begin{array}{ll}
r(A)<r\big(A\mid\boldsymbol b\big) & \text{无解}\\[3pt]
r(A)=r\big(A\mid\boldsymbol b\big)=n & \text{唯一解}\\[3pt]
r(A)=r\big(A\mid\boldsymbol b\big)<n & \text{有无穷多解}
\end{array}}
\qquad\qquad
\bx{r\big(A\mid\boldsymbol b\big)-r(A)\le1}
\]
\[
r\big(A\mid\boldsymbol b\big)=\max\{r(A),r\big(A\mid\boldsymbol b\big)\}
\]
\[
\bx{
\text{通解}\ x=\boldsymbol x_{0}
+k_{1}\boldsymbol\eta_{1}+\cdots+k_{n-r(A)}\boldsymbol\eta_{n-r(A)}
}
\quad
\begin{gathered}
\boldsymbol x_{0}:\ Ax_{0}=\boldsymbol b\ \text{的一个特解}\\
\boldsymbol\eta_{1},\dots,\boldsymbol\eta_{n-r(A)}:\ Ax=\boldsymbol 0\ \text{的一组基础解系}\\
\text{自由度：}\ n-r(A)
\end{gathered}
\]
\[
\text{解的性质：}\quad
Ax=\boldsymbol b\ \text{通解}\
\boldsymbol x_{1},\boldsymbol x_{2}
\ \Rightarrow\ \text{它们的任意线性组合仍为解；}\quad
k(\boldsymbol x_{1}-\boldsymbol x_{2})\ \text{为}\ Ax=\boldsymbol 0\ \text{的解}
\]
\end{formbox}

\begin{formbox}
\hd{含参数方程组的讨论（通法）}
\[
\left(\begin{array}{c|c}A&\boldsymbol b\end{array}\right)
\ \overset{\text{初等行变换}}{\longrightarrow}\ \text{阶梯形}
\quad\text{再比较}\quad r(A)\ \text{与}\ r\big(A\mid\boldsymbol b\big)
\]
\[
\left\{
\begin{array}{ll}
\text{出现}\ 0=1\ \text{或}\ 0=k\ (k\ne0) & \Rightarrow\ \text{无解}\\[4pt]
\text{无矛盾行，}\ r(A)=r\big(A\mid\boldsymbol b\big)=n & \Rightarrow\ \text{唯一解}\\[4pt]
\text{无矛盾行，}\ r(A)=r\big(A\mid\boldsymbol b\big)<n & \Rightarrow\ \text{无穷多解}
\end{array}
\right.
\]
\end{formbox}
\end{secbox}

% -------------------------------------------------------------------
\begin{secbox}[十五、线性代数 · \textcircled{5} 特征值与特征向量]

\begin{formbox}
\[
\exists\,\lambda,\ \boldsymbol\xi\ne\boldsymbol 0,\quad
\bx{A\boldsymbol\xi=\lambda\boldsymbol\xi}
\quad\Longleftrightarrow\quad
|\lambda E-A|=0
\]
$\lambda$ 为 $A$ 的特征值，$\boldsymbol\xi$ 为属于 $\lambda$ 的特征向量。
\[
\text{特征多项式}\ f(\lambda)=|\lambda E-A|
=\lambda^{n}-a_{11}\lambda^{n-1}+\cdots+(-1)^{n}|A|
\]
\[
\text{特征值}\ \lambda_{1},\dots,\lambda_{n}\text{ 的基本关系：}\quad
\bx{\sum_{i=1}^{n}\lambda_{i}=\sum_{i=1}^{n}a_{ii}=\operatorname{tr}(A)
=\operatorname{tr}(A^{k})\ \ (k\in\mathbb N^{*})}
\qquad
\bx{\prod_{i=1}^{n}\lambda_{i}=|A|}
\]
\[
\lambda_{1}+\cdots+\lambda_{n}=
\ \lambda_{1}^{k}+\cdots+\lambda_{n}^{k}
=\operatorname{tr}(A^{k})
\qquad\text{（}\ k\ \text{重特征值按重数计）}
\]
\[
\begin{aligned}
&|A|\ne0\ \Rightarrow\ \frac1{\lambda}\ \text{是}\ A^{-1}\ \text{的特征值}\\[3pt]
&\qquad \frac{|A|}{\lambda}\ \text{是}\ A^{*}\ \text{的特征值}\\[3pt]
&\lambda^{k}\ \text{是}\ A^{k}\ \text{的特征值}\\[3pt]
&\lambda_{1}\ne\lambda_{2}\ \Rightarrow\
\xi_{1},\xi_{2}\ \text{分别对应的特征向量}\ \xi_{1}^{T}\xi_{2}=0\\[3pt]
&\xi,\eta\ \text{为}\ \lambda\ \text{的特征向量}\ \Rightarrow\
&k\xi+l\eta\ \text{为}\ \lambda\ \text{的特征向量（当两者线性无关时）}\\[3pt]
&\xi\ \text{为}\ A\ \text{的特征向量}\ \Rightarrow\ \xi\ \text{为}\ A^{T}\ \text{的特征向量}\\[3pt]
&\text{上（下）三角矩阵的特征值：}\ \text{主对角线元素}
\end{aligned}
\]
\end{formbox}

\begin{formbox}
\hd{相似矩阵与相似对角化}
\[
B=P^{-1}AP\iff PA=BP
\quad\Longrightarrow\quad P\ \text{可逆}
\qquad\text{称}\ A\ \text{与}\ B\ \text{相似，记}\ A\sim B
\]
\[
A\sim B\ \Rightarrow\
\begin{array}{l}
r(A)=r(B)\\
|A|=|B|,\quad \operatorname{tr}(A)=\operatorname{tr}(B)\\
\operatorname{tr}(A^{k})=\operatorname{tr}(B^{k})\\
\text{特征值相同（含重数）}\\
A^{k}\sim B^{k}
\end{array}
\qquad\qquad
\bx{B=P^{-1}AP\ \Rightarrow\ A^{k}=P\,B^{k}\,P^{-1}}
\]
\[
\bx{
A\ \text{可相似对角化}
\iff \exists\ n\ \text{个线性无关的特征向量}
}
\]
\[
\bx{k\ \text{重特征值}\ \lambda\ \text{有}\ k\ \text{个线性无关的特征向量}}
\]
\[
\bx{
A=P\Lambda P^{-1},\quad
\Lambda=\operatorname{diag}(\lambda_{1},\dots,\lambda_{n}),\ P\ \text{可逆}
}
\qquad
A^{k}=P\Lambda^{k}P^{-1},
\quad
\Lambda^{k}=\operatorname{diag}(\lambda_{1}^{k},\dots,\lambda_{n}^{k})
\]
\[
P_{1},\Lambda_{1}\ \text{与}\ P_{2},\Lambda_{2}\ \text{分别把}\ A\ \text{对角化}
\ \Rightarrow\ P_{2}^{-1}P_{1}\ \text{与}\ \Lambda_{2}^{-1}\Lambda_{1}\ \text{均可交换}
\]
\hd{相似与合同的区别（高频易混）}
\begin{center}
\renewcommand{\arraystretch}{1.5}
\begin{tabular}{|c|c|c|}
\hline
 & 相似 $A\sim B$ & 合同 \\ \hline
关系式 & $B=P^{-1}AP$ & $B=P^{T}AP$ \\ \hline
允许的变换 & 初等行变换 & 初等行\textbf{列}变换 \\ \hline
保持不变 & 特征值、行列式、迹、秩 & 秩、行列式、迹、标准形系数 \\ \hline
典型用途 & 相似对角化、求 $A^{k}$ & 化二次型为标准形 \\ \hline
$P$ 为正交矩阵时 & 正交相似对角化 & 同时也是合同变换 \\ \hline
\end{tabular}
\end{center}
\end{formbox}

\begin{formbox}
\hd{实对称矩阵 $A$（$A^{T}=A$）的性质}
\[
\begin{gathered}
\boxed{\text{特征值全为实数}\qquad
\text{不同特征值对应的特征向量相互正交}\qquad
\text{一定可相似对角化}}\\[5pt]
\text{（可用正交矩阵}\ Q:/ Q^{T}AQ=\Lambda\text{）}\\[5pt]
\lambda_{1}\ne\lambda_{2}\ \Rightarrow\ \xi_{1}^{T}\xi_{2}=0\\[4pt]
k\ \text{重特征值对应}\ k\ \text{维特征子空间，可从中取出}\ k\ \text{个正交单位特征向量}\\[4pt]
|A|\ \text{为实数};\qquad |\lambda E-A|\ \text{为实系数多项式};\qquad
Q^{T}AQ=\Lambda\ \Rightarrow\ A\ \text{可合同对角化}
\end{gathered}
\]
\[
\text{重根情形：若}\ \xi_{1},\xi_{2}\ \text{不正交，需先施密特正交化再单位化}
\]
\end{formbox}
\end{secbox}

% -------------------------------------------------------------------
\begin{secbox}[十六、线性代数 · \textcircled{6} 二次型]

\begin{formbox}
\[
f(x_{1},\dots,x_{n})=\sum_{i=1}^{n}a_{ii}x_{i}^{2}
+2\sum_{i<j}a_{ij}x_{i}x_{j}
\quad\Longleftrightarrow\quad
\bx{f(\boldsymbol x)=\boldsymbol x^{T}A\boldsymbol x},
\qquad A=A^{T}
\]
\[
f=\boldsymbol x^{T}A\boldsymbol x,\ \ \boldsymbol y^{T}B\boldsymbol y
\quad(A,B\ \text{对称})\ \Rightarrow\
A\ \text{与}\ B\ \text{合同}\ \Longrightarrow\ f\equiv g
\quad\Longleftrightarrow\quad \boldsymbol x=C\boldsymbol y,\ A=C^{T}BC
\]
\[
A\ \text{对称}\iff\ \text{二次型矩阵唯一确定}:
\quad f\ \text{写成}\ \boldsymbol x^{T}A\boldsymbol x,\ A\ \text{可取}\ \frac{A+A^{T}}{2}
\]
\[
r(f)=r(A)=\text{非零系数项的个数}=\text{标准形中非零系数个数}
=\text{规范形中系数为}\ \pm1\ \text{的个数}
\]
\end{formbox}

\begin{formbox}
\hd{标准形、规范形与惯性定理}
\[
\begin{gathered}
\text{标准形：}\quad f=\lambda_{1}y_{1}^{2}+\cdots+\lambda_{r}y_{r}^{2}\ \ (\lambda_{i}\ne0)\\[4pt]
\text{规范形：}\quad f=y_{1}^{2}+\cdots+y_{p}^{2}-y_{p+1}^{2}-\cdots-y_{r}^{2}\\[4pt]
\bx{\text{惯性定理：}\quad r,\ p,\ q\ \text{在合同变换下均保持不变}\ (r=p+q)}
\end{gathered}
\]
\[
\bx{
\text{任意合同变换}\ P^{T}AP=B\ \Rightarrow\ r(A)=r(B)
}
\]
\[
\bx{\text{正惯性指数相同，负惯性指数相同}}
\]
\[
\text{实对称矩阵的任一特征向量系（含重根的}\ k\ \text{维正交单位基）}
\ \Longrightarrow \ \text{正交相似对角化}
\]
\end{formbox}

\begin{formbox}
\hd{正交变换化二次型为标准形（数三重点流程）}
\[
\bxe{
\begin{gathered}
\text{① 写出二次型矩阵}\ A\ (A=A^{T})\\[4pt]
\text{② 求特征值：}\ |\lambda E-A|=0\ \Rightarrow\ \lambda_{1},\lambda_{2},\lambda_{3}
\quad\textcolor{black!55}{\textstyle \sum\lambda_{i}=\operatorname{tr}(A),\ \prod\lambda_{i}=|A|}\\[4pt]
\text{③ 求特征向量：}\ (\lambda_{i}E-A)\boldsymbol\xi_{i}=\boldsymbol 0\\[4pt]
\text{④ 重根时}\ \textbf{施密特正交化};\qquad
\text{⑤ 单位化得}\ \boldsymbol\eta_{1},\boldsymbol\eta_{2},\boldsymbol\eta_{3}\\[4pt]
\text{⑥ } Q=(\boldsymbol\eta_{1}\,\boldsymbol\eta_{2}\,\boldsymbol\eta_{3})\ \text{为正交矩阵，}\ 
Q^{T}AQ=\Lambda=\operatorname{diag}(\lambda_{1},\lambda_{2},\lambda_{3})\\[4pt]
\Rightarrow\ \boldsymbol x=Q\boldsymbol y,\qquad
\bx{f=\lambda_{1}y_{1}^{2}+\lambda_{2}y_{2}^{2}+\lambda_{3}y_{3}^{2}}
\end{gathered}}
\]
\[
Q^{T}Q=E
\quad\Rightarrow\quad
|\boldsymbol x_{1}|^{2}+|\boldsymbol x_{2}|^{2}+|\boldsymbol x_{3}|^{2}
=|\boldsymbol y_{1}|^{2}+|\boldsymbol y_{2}|^{2}+|\boldsymbol y_{3}|^{2}
\quad\text{（}\ Q\ \text{保持长度与夹角）}
\]
\hd{配方法化标准形}
\[
\bx{\text{配方（变量平方项配方）}\ \longrightarrow\ \text{变量非平方项配方（凑完全平方）}
\ \longrightarrow\ \text{标准形}}
\]
\[
\boldsymbol x=C\boldsymbol y\ (C\ \text{可逆})\ \Rightarrow\ C^{T}AC=\Lambda\ (\text{标准形}),
\qquad \text{注意 } C\ \text{是合同变换矩阵，}\ \text{不是正交矩阵}
\]
\end{formbox}

\begin{formbox}
\hd{正定与半正定（$A$ 为实对称矩阵）}
\[
\bx{
A\ \text{正定}
\iff \forall\,\boldsymbol x\ne\boldsymbol 0\,\ \boldsymbol x^{T}A\boldsymbol x>0
\iff \lambda_{1},\dots,\lambda_{n}>0
}
\]
\[
\bx{
\iff r(A)=n\ \text{且}\ |A|>0
\iff \Delta_{k}>0,\ k=1,\dots,n
\iff A=C^{T}C\ (C\ \text{可逆})
}
\]
\[
\bx{
A\ \text{正定}\ \Rightarrow
a_{ii}>0,\ \boldsymbol x^{T}A\boldsymbol x\ \text{无零点},\ \lambda_{max}>0
}
\]
\[
\bx{|A_{i}|>0,\ \text{即任意阶顺序主子式均正}}
\]
\[
\begin{aligned}
&A\ \text{半正定}
\iff \boldsymbol x^{T}A\boldsymbol x\ge0,\ \forall\boldsymbol x
\iff \lambda_{i}\ge0,\ \forall i
\iff A=C^{T}C
\iff A\ \text{的所有}\ \textbf{主子式}\ \ge0\\[4pt]
&A\ \text{正定}\iff A\ \text{半正定且}\ |A|>0;
\qquad A\ \text{正定}\Rightarrow A\ \text{半正定}\\[4pt]
&A\ \text{负定}\iff (-1)^{k}\Delta_{k}>0,\ k=1,\dots,n
\iff \lambda_{i}<0,\ \forall i\\[4pt]
&A\ \text{正定}+A\ \text{正定}=A\ \text{正定};\qquad
A\ \text{正定},\ k>0\ \Rightarrow\ kA\ \text{正定};\qquad
A\ \text{正定}\ \Rightarrow\ A^{-1}\ \text{正定}\\[4pt]
&A\ \text{正定},\ B\ \text{任意实对称}\ \Rightarrow\ A\pm B\ \text{不确定}
\quad\Rightarrow\ \bx{\text{不能把一般矩阵与正定矩阵随意加减}}
\end{aligned}
\]
\end{formbox}
\end{secbox}
